\documentclass[11pt]{article}

\usepackage[margin=1in]{geometry}
\usepackage{amsmath,amssymb,amsthm}
\usepackage{enumitem}
\usepackage{hyperref}
\usepackage{xcolor}
\usepackage{booktabs}

\hypersetup{colorlinks=true, linkcolor=blue, urlcolor=blue, citecolor=blue}

\title{TOAE209/TOAH209 -- Design and Analysis of Algorithms\\[4pt]
\large Week 12: Theoretical Questions}
\author{Foreign Trade University}
\date{}

\begin{document}
\maketitle

\section*{Instructions}

Part A contains 17 questions requiring written explanations or short proofs. Part B is a
10-question quiz (true/false and multiple choice). Attempt every question on your own
before consulting Part C (Solutions) at the end of this document.

\section*{Part A: Theory Questions}

\begin{enumerate}[label=\textbf{A\arabic*.}, leftmargin=2.2em]

\item Define a \emph{polynomial-time reduction} $X \le_p Y$. State precisely the two
consequences of $X \le_p Y$ regarding the tractability of $X$ and $Y$, and explain which
of the two is used to \emph{prove a problem hard}.

\item Explain why the theory of NP-completeness restricts attention to \emph{decision
problems} (yes/no answers). Show how the optimization problem ``find a maximum independent
set'' can be solved using a polynomial number of calls to the decision problem ``does $G$
have an independent set of size $\ge k$?''.

\item Define the class $\mathrm{P}$ and the class $\mathrm{NP}$ in terms of algorithms,
certificates, and verifiers. In one sentence each, contrast what $\mathrm{P}$ and
$\mathrm{NP}$ capture.

\item Prove that $\mathrm{P} \subseteq \mathrm{NP}$.

\item For the \textsc{Independent Set} problem, describe explicitly: (a) what a certificate
is, (b) what the verifier checks, and (c) why the verifier runs in polynomial time. Do the
same for \textsc{3-SAT}.

\item State precisely what it means for a problem to be \emph{NP-hard} and what it means to
be \emph{NP-complete}. Why does NP-completeness require a separate proof that the problem is
\emph{in} $\mathrm{NP}$, on top of the hardness argument?

\item State the Cook--Levin theorem. In two or three sentences, describe the high-level
idea of why \emph{every} problem in $\mathrm{NP}$ reduces to \textsc{SAT}.

\item Prove that $\le_p$ is transitive: if $X \le_p Y$ and $Y \le_p Z$ then $X \le_p Z$.
Then explain how transitivity, together with one known NP-complete problem, makes
subsequent NP-completeness proofs much shorter.

\item State the standard ``recipe'' (the list of steps) for proving that a problem $Y$ is
NP-complete. Identify which step establishes membership in $\mathrm{NP}$ and which
establishes NP-hardness.

\item Prove the complementarity lemma: in any graph $G = (V, E)$ with $|V| = n$, a set $S$
is an independent set iff $V \setminus S$ is a vertex cover. Deduce that $G$ has an
independent set of size $\ge k$ iff it has a vertex cover of size $\le n - k$, and hence
that \textsc{Independent Set} $\le_p$ \textsc{Vertex Cover}.

\item Prove that $S$ is an independent set in $G$ iff $S$ is a clique in the complement
graph $\overline{G}$, and conclude that \textsc{Independent Set} $\le_p$ \textsc{Clique}.

\item Describe the reduction \textsc{3-SAT} $\le_p$ \textsc{Independent Set}: how many
vertices does the constructed graph have, what are the ``triangle'' edges, what are the
``conflict'' edges, and what is the target size $k$? Prove the forward direction (a
satisfying assignment yields an independent set of size $k$).

\item For the same reduction, prove the reverse direction: an independent set of size $k$
yields a satisfying assignment. In particular, explain why such an independent set must
contain exactly one vertex from each clause-triangle, and why the conflict edges guarantee
the resulting assignment is consistent.

\item Sketch the reduction \textsc{Hamiltonian Cycle} $\le_p$ \textsc{TSP} (decision
version). What distances do you assign to edges and non-edges of the graph, what tour-length
bound $D$ do you use, and why does a tour of length $\le D$ correspond exactly to a
Hamiltonian cycle?

\item \textsc{Subset Sum} has a dynamic-programming algorithm running in time $O(nW)$, where
$W$ is the target. Explain why this does \emph{not} contradict the NP-completeness of
\textsc{Subset Sum}. Define \emph{pseudo-polynomial} time and explain the distinction
between the numeric value $W$ and the input size.

\item A student writes: ``To prove \textsc{Vertex Cover} is NP-hard, I reduced
\textsc{Vertex Cover} to \textsc{3-SAT}; since \textsc{3-SAT} is NP-complete, so is
\textsc{Vertex Cover}.'' Identify the error precisely, and state what reduction the student
should have constructed instead.

\item \emph{Discussion (open-ended).} You are handed a brand-new optimization problem at
work and asked whether there is an efficient algorithm. Describe a systematic process for
(a) deciding whether to search for a polynomial-time algorithm or for an NP-completeness
proof, and (b) if you suspect hardness, how you would go about constructing a reduction.
What known NP-complete problems would you try first for a graph problem? a packing/number
problem? an ordering/routing problem?

\end{enumerate}

\newpage
\section*{Part B: Quiz}

\begin{enumerate}[label=\textbf{B\arabic*.}, leftmargin=2.2em]

\item \textbf{(True/False)} If $X \le_p Y$ and $Y$ has a polynomial-time algorithm, then
$X$ has a polynomial-time algorithm.

\item \textbf{(Multiple choice)} The class $\mathrm{NP}$ is the set of decision problems
for which:
\begin{enumerate}[label=(\alph*)]
    \item a solution can be found in polynomial time
    \item a proposed solution (certificate) can be verified in polynomial time
    \item no polynomial-time algorithm can possibly exist
    \item the answer is always ``yes''
\end{enumerate}

\item \textbf{(True/False)} Proving $\mathrm{P} = \mathrm{NP}$ would follow from finding a
single polynomial-time algorithm for any one NP-complete problem.

\item \textbf{(Multiple choice)} The Cook--Levin theorem establishes that:
\begin{enumerate}[label=(\alph*)]
    \item every problem in $\mathrm{NP}$ is in $\mathrm{P}$
    \item \textsc{SAT} (and \textsc{3-SAT}) is NP-complete
    \item \textsc{Subset Sum} can be solved in polynomial time
    \item $\mathrm{P} \ne \mathrm{NP}$
\end{enumerate}

\item \textbf{(True/False)} To prove a problem $Y$ is NP-hard, it suffices to give a
polynomial-time reduction \emph{from} a known NP-complete problem \emph{to} $Y$.

\item \textbf{(Short answer)} In a graph $G$ on $n$ vertices, if the maximum independent
set has size $13$, what is the size of the minimum vertex cover (in terms of $n$)?

\item \textbf{(Multiple choice)} In the reduction \textsc{3-SAT} $\le_p$ \textsc{Independent
Set}, a 3-SAT formula with $k$ clauses produces a graph in which we look for an independent
set of size:
\begin{enumerate}[label=(\alph*)]
    \item $1$
    \item $3$
    \item $k$
    \item $3k$
\end{enumerate}

\item \textbf{(True/False)} A set $S$ is a clique in $G$ if and only if $S$ is an
independent set in the complement graph $\overline{G}$.

\item \textbf{(Short answer)} \textsc{2-Coloring} is in $\mathrm{P}$ but \textsc{3-Coloring}
is NP-complete. What classical polynomial-time test decides 2-colorability?

\item \textbf{(Multiple choice)} An $O(nW)$ algorithm for \textsc{Subset Sum} (with target
$W$) is best described as:
\begin{enumerate}[label=(\alph*)]
    \item polynomial in the input size
    \item pseudo-polynomial (polynomial in the numeric value $W$, not in $\log W$)
    \item proof that $\mathrm{P} = \mathrm{NP}$
    \item exponential in $n$ for all inputs
\end{enumerate}

\end{enumerate}

\newpage
\section*{Part C: Solutions}

\subsection*{Part A}

\begin{enumerate}[label=\textbf{A\arabic*.}, leftmargin=2.2em]

\item $X \le_p Y$ means an arbitrary instance of $X$ can be solved using a polynomial
number of basic steps plus a polynomial number of calls to an oracle (subroutine) for $Y$.
The two consequences: (a) if $Y \in \mathrm{P}$ then $X \in \mathrm{P}$ (tractability flows
forward); (b) if $X \notin \mathrm{P}$ then $Y \notin \mathrm{P}$ (intractability flows
backward). To \emph{prove a problem $Y$ hard}, we use (b): we take a known-hard $X$ and
show $X \le_p Y$, so $Y$ is at least as hard as $X$.

\item Decision problems have a single bit of output, which makes ``polynomial time'' clean
to define and lets reductions transform yes/no answers into yes/no answers. The restriction
loses nothing essential because optimization reduces to decision: to find the size of a
maximum independent set in $G$ on $n$ vertices, ask the decision oracle ``is there an
independent set of size $\ge k$?'' for $k = n, n-1, \ldots$ (or binary-search on $k$ in
$O(\log n)$ queries) to find the largest feasible $k^*$. To recover an actual set, fix
vertices one at a time: for each vertex $v$, delete it and ask whether the remaining graph
still has an independent set of size $k^*$; keep $v$ out (it was not needed) if so,
otherwise it belongs to the set -- a polynomial number of oracle calls in total.

\item $\mathrm{P}$ is the set of decision problems solvable by an algorithm that always
halts in time polynomial in the input length. $\mathrm{NP}$ is the set of decision problems
that have an \emph{efficient verifier}: a polynomial-time algorithm $B(s, t)$ such that $s$
is a yes-instance iff there exists a polynomial-length certificate $t$ with $B(s,t) =
\mathtt{yes}$. In one sentence each: $\mathrm{P}$ captures problems we can efficiently
\emph{solve}; $\mathrm{NP}$ captures problems whose solutions we can efficiently
\emph{check} once handed a candidate.

\item Let $X \in \mathrm{P}$ via a polynomial-time algorithm $A$. Define a verifier $B(s,t)
:= A(s)$ that ignores the certificate $t$ (take $t$ to be the empty string). Then $B$ runs
in polynomial time, and $X(s) = \mathtt{yes}$ iff there exists a certificate (trivially,
the empty one) with $B(s,t) = \mathtt{yes}$, namely exactly when $A(s) = \mathtt{yes}$.
Hence $X \in \mathrm{NP}$, so $\mathrm{P} \subseteq \mathrm{NP}$.

\item \textsc{Independent Set}, input $(G, k)$: (a) the certificate is a set $S$ of
vertices; (b) the verifier checks $|S| \ge k$ and that no edge of $G$ has both endpoints in
$S$; (c) this takes $O(|S|^2)$ or $O(|E|)$ time -- polynomial. \textsc{3-SAT}, input a CNF
formula $\Phi$: (a) the certificate is a truth assignment to the variables; (b) the verifier
evaluates each clause under the assignment and checks every clause is satisfied; (c) this is
linear in the size of $\Phi$. In both cases the certificate has length polynomial in the
input and the check is polynomial, so both problems are in $\mathrm{NP}$.

\item $Y$ is \emph{NP-hard} if $X \le_p Y$ for every $X \in \mathrm{NP}$ (it is at least as
hard as everything in $\mathrm{NP}$). $Y$ is \emph{NP-complete} if it is NP-hard \emph{and}
$Y \in \mathrm{NP}$. We separately prove membership in $\mathrm{NP}$ because NP-hardness
alone does not place $Y$ inside $\mathrm{NP}$ -- there are NP-hard problems (e.g.\ certain
undecidable or PSPACE-hard problems) that are not in $\mathrm{NP}$ at all. ``Complete''
means hardest \emph{within} the class, which requires being a member of the class.

\item \textbf{Cook--Levin theorem:} \textsc{SAT} (and \textsc{3-SAT}) is NP-complete.
High-level idea: take any $X \in \mathrm{NP}$ with efficient verifier $B$ running in
polynomial time $p(|s|)$. The entire computation of $B$ on input $s$ and certificate $t$ --
the memory contents at each of the $p(|s|)$ time steps -- can be encoded with polynomially
many Boolean variables, and a polynomial-size CNF formula can express ``these variables
describe a valid accepting computation of $B$ for some certificate $t$.'' That formula is
satisfiable iff a valid certificate exists iff $s$ is a yes-instance, giving $X \le_p
\textsc{SAT}$.

\item \textbf{Transitivity.} Suppose $X \le_p Y$ and $Y \le_p Z$. Solve an instance of $X$
using the first reduction: polynomially many steps and polynomially many calls to a
$Y$-oracle. Replace each $Y$-oracle call by the second reduction, which solves $Y$ using
polynomially many steps and polynomially many calls to a $Z$-oracle. Composing, $X$ is
solved with a polynomial number of $Z$-oracle calls and polynomial overhead, so $X \le_p
Z$. \emph{Why it shortens proofs:} once a single problem (\textsc{SAT}, via Cook--Levin) is
known to be NP-complete, to show a new $Y$ is NP-hard we need only exhibit \emph{one}
reduction from \emph{one} known NP-complete problem $X$ to $Y$; transitivity then chains it
back through $X$ to \emph{every} problem in $\mathrm{NP}$, so we never have to re-encode
arbitrary NP computations.

\item \textbf{Recipe to prove $Y$ NP-complete:} (1) state $Y$ as a decision problem; (2)
prove $Y \in \mathrm{NP}$ by giving a short certificate and a polynomial-time verifier
[\emph{this is the membership step}]; (3) pick a known NP-complete problem $X$; (4)
construct a polynomial-time reduction $X \le_p Y$ (the gadget); (5) prove correctness in
both directions ($x$ yes $\iff f(x)$ yes); (6) check the reduction is polynomial-time
[\emph{steps 3--6 establish NP-hardness}]. Membership (step 2) plus NP-hardness (steps 3--6)
together give NP-completeness.

\item \textbf{Complementarity.} ($\Rightarrow$) Let $S$ be independent. For any edge
$(u,v) \in E$, $u$ and $v$ are not both in $S$ (independence), so at least one is in $V
\setminus S$ -- hence $V \setminus S$ covers every edge, i.e.\ it is a vertex cover.
($\Leftarrow$) Let $V \setminus S$ be a vertex cover. If some edge had both endpoints in
$S$, that edge would be uncovered, a contradiction; so no edge lies inside $S$, i.e.\ $S$
is independent. Since $|V \setminus S| = n - |S|$, an independent set of size $\ge k$
corresponds exactly to a vertex cover of size $\le n - k$. This gives the reduction
\textsc{Independent Set} $\le_p$ \textsc{Vertex Cover}: map $(G, k) \mapsto (G, n - k)$
(constant work plus one oracle call), preserving the answer.

\item $S$ is independent in $G$ $\iff$ no pair of vertices in $S$ is an edge of $G$ $\iff$
every pair of vertices in $S$ is a \emph{non-edge} of $G$ $\iff$ every pair in $S$ is an
edge of $\overline{G}$ (whose edges are precisely the non-edges of $G$) $\iff$ $S$ is a
clique in $\overline{G}$. The reduction \textsc{Independent Set} $\le_p$ \textsc{Clique}
maps $(G, k) \mapsto (\overline{G}, k)$; building $\overline{G}$ takes $O(n^2)$ time, and a
$k$-independent set in $G$ exists iff a $k$-clique in $\overline{G}$ exists.

\item \textbf{Construction.} For a formula with $k$ clauses, build a graph with $3k$
vertices: three vertices per clause, one per literal. \emph{Triangle edges}: within each
clause, join all three of its literal-vertices pairwise (forming a triangle). \emph{Conflict
edges}: between any two vertices in different clauses whose literals are complementary (one
is $x_i$, the other $\overline{x_i}$), add an edge. Target size $k$ = number of clauses.

\textbf{Forward direction.} Suppose $\Phi$ has a satisfying assignment. Each clause has at
least one literal evaluating to true; pick one such literal per clause and take its vertex.
This selects exactly one vertex from each triangle (so no triangle edge is violated, since
triangle edges only connect vertices \emph{within} a clause and we picked one), and no two
selected vertices conflict: a conflict edge connects $x_i$ and $\overline{x_i}$, but under a
single assignment these literals cannot both be true, so we never selected both. Hence the
selected $k$ vertices form an independent set of size $k$.

\item \textbf{Reverse direction.} Suppose $G$ has an independent set $S$ with $|S| = k$.
Within each clause-triangle the three vertices are mutually adjacent, so $S$ contains at
most one vertex per triangle. There are exactly $k$ triangles and $|S| = k$, so $S$ contains
\emph{exactly one} vertex per triangle. Construct an assignment by setting, for each vertex
in $S$ labelled with literal $\ell$, the underlying variable so that $\ell$ is true (e.g.\
if $\ell = \overline{x_i}$ set $x_i$ false). This is well-defined and consistent: if $S$
contained both an $x_i$-vertex and an $\overline{x_i}$-vertex, the conflict edge between
them would violate independence -- but $S$ is independent, so no such pair exists, and we
never demand contradictory values of any variable. Unconstrained variables are set
arbitrarily. Since $S$ has a vertex in every clause-triangle and that vertex's literal is
made true, every clause is satisfied, so $\Phi$ is satisfiable.

\item \textbf{Reduction \textsc{Hamiltonian Cycle} $\le_p$ \textsc{TSP}.} Given a graph $G
= (V, E)$ with $|V| = n$, create one city per vertex. Set the distance $d(i, j) = 1$ if
$(i, j) \in E$ and $d(i, j) = 2$ (any value $> 1$, e.g.\ $\infty$) if $(i, j) \notin E$.
Set the bound $D = n$. A tour visits all $n$ cities and returns to start, using exactly $n$
edges; its length is $\ge n$, with equality iff \emph{all} $n$ tour-edges have distance $1$,
i.e.\ all are edges of $G$. So a tour of length $\le n$ exists iff there is a cycle through
all vertices using only edges of $G$ -- precisely a Hamiltonian cycle. The construction is
$O(n^2)$, hence polynomial, and the equivalence is exact.

\item The $O(nW)$ algorithm is \emph{pseudo-polynomial}: it is polynomial in $n$ and in the
numeric \emph{value} $W$, but the \emph{input size} is the number of bits used to write the
numbers, roughly $\sum_i \log w_i + \log W$, i.e.\ proportional to $\log W$, not $W$. When
$W \approx 2^{m}$ for input size $\approx m$ bits, $O(nW) = O(n \cdot 2^m)$ is exponential
in the input size. So the DP is fast only when the numbers are small (written in unary, or
bounded by a polynomial in $n$); for the general problem, with numbers given in binary, it
is exponential. This is fully consistent with NP-completeness, which is about polynomiality
in the input \emph{size}.

\item The error is the \emph{direction} of the reduction. The student reduced \textsc{Vertex
Cover} $\le_p$ \textsc{3-SAT}, which only shows \textsc{Vertex Cover} is no \emph{harder}
than \textsc{3-SAT} (i.e.\ \textsc{Vertex Cover} $\in \mathrm{NP}$, which we already knew);
it says nothing about \textsc{Vertex Cover} being hard. To prove \textsc{Vertex Cover}
NP-hard, the student must reduce a known NP-complete problem \emph{to} \textsc{Vertex
Cover}, e.g.\ \textsc{3-SAT} $\le_p$ \textsc{Vertex Cover} (or \textsc{Independent Set}
$\le_p$ \textsc{Vertex Cover} via complementarity). Hardness flows \emph{from} the known
problem \emph{into} the target, so the known-hard problem must be the source of the arrow.

\item A systematic process. \textbf{(a) Decide which way to bet.} First nail down the exact
decision version and the input encoding. Try small instances by hand and look for structure
that matches a paradigm you know (greedy, DP, flow, matroid) -- if a natural subproblem
structure or exchange argument appears, pursue a polynomial algorithm. Simultaneously, look
for ``combinatorial explosion'' signals: a need to choose a subset/permutation/assignment
with global interacting constraints, the appearance of words like ``exactly,'' ``cover,''
``Hamiltonian,'' ``3-coloring,'' or an obvious resemblance to a known NP-complete problem.
If polynomial attempts keep failing and the problem smells like SAT/packing/routing, switch
to seeking an NP-completeness proof. \textbf{(b) Build the reduction.} Identify the source
problem whose ``logic'' your problem can encode, design gadgets that translate that
problem's variables/constraints into your problem's objects, and verify both directions.
\textbf{First problems to try:} for a \emph{graph} problem -- \textsc{Independent Set},
\textsc{Vertex Cover}, \textsc{Clique}, \textsc{3-Coloring}, or \textsc{Hamiltonian Cycle};
for a \emph{packing/number} problem -- \textsc{Subset Sum} or \textsc{Knapsack}; for an
\emph{ordering/routing} problem -- \textsc{Hamiltonian Cycle} or \textsc{TSP}; and
\textsc{3-SAT} as the universal fallback when the structure is logical/constraint-based.

\end{enumerate}

\subsection*{Part B}

\begin{enumerate}[label=\textbf{B\arabic*.}, leftmargin=2.2em]

\item \textbf{True.} This is the forward direction of a reduction: tractability of $Y$
transfers to $X$ (Proposition: $X \le_p Y$ and $Y \in \mathrm{P}$ imply $X \in \mathrm{P}$).

\item \textbf{(b)} A proposed solution (certificate) can be verified in polynomial time.
This is the verifier-based definition of $\mathrm{NP}$; it does \emph{not} require that
solutions can be \emph{found} quickly.

\item \textbf{True.} If any NP-complete problem is in $\mathrm{P}$, then since every problem
in $\mathrm{NP}$ reduces to it, all of $\mathrm{NP}$ is in $\mathrm{P}$, so $\mathrm{P} =
\mathrm{NP}$.

\item \textbf{(b)} \textsc{SAT} (and \textsc{3-SAT}) is NP-complete. Cook--Levin provides
the first NP-complete problem; it does not resolve $\mathrm{P}$ vs.\ $\mathrm{NP}$.

\item \textbf{True.} NP-hardness of $Y$ is exactly: some known NP-complete $X$ satisfies $X
\le_p Y$ (and by transitivity all of $\mathrm{NP}$ reduces to $Y$). The arrow goes
\emph{from} the known-hard problem \emph{to} $Y$.

\item The minimum vertex cover has size $n - 13$, by complementarity ($|$min VC$| = n -
|$max IS$|$).

\item \textbf{(c) $k$.} The constructed graph has $3k$ vertices (a triangle per clause), but
the target independent-set size is $k$ -- one vertex per clause.

\item \textbf{True.} A clique in $G$ is a set in which all pairs are edges of $G$, i.e.\ no
pair is an edge of $\overline{G}$, i.e.\ an independent set in $\overline{G}$ (and vice
versa).

\item Check whether the graph is \textbf{bipartite} (no odd cycle), via BFS/DFS 2-coloring.
A graph is 2-colorable iff it is bipartite, which is decidable in linear time.

\item \textbf{(b)} Pseudo-polynomial: polynomial in the numeric value $W$ but exponential
in the input size $\log W$ when $W$ is given in binary -- consistent with \textsc{Subset
Sum} being NP-complete.

\end{enumerate}

\end{document}
